\section{Diagram 1}\label{sec:d1}

Diagram~1 is one of the two disconnected diagrams left by \cite[\S7.1, p.~565]{AK}: the
$z$-stroke $12$ with $z$-circles at $1$ and $2$, and the $w$-triangle with strokes $34,35,45$ and
no circle; there are no other strokes or circles.
We keep the numbering of \cite[\S8.1]{AK}: $P=\{1,2\}$, $T=\{3,4,5\}$, the letter $i$ denotes a
body of $P$ and $k,l$ denote bodies of $T$.  Put $M_P=m_1+m_2$, $M_T=m_3+m_4+m_5$,
$Z_P=M_P^{-1}(m_1z_1+m_2z_2)$, $Z_T=M_T^{-1}\sum_km_kz_k$, and define $W_P,W_T$ likewise.  On
$A_m$ we have $\sum_jm_jz_j=\sum_jm_jw_j=0$ (proof of Lemma~\ref{lem:distances}), so
$M_PZ_P+M_TZ_T=0=M_PW_P+M_TW_T$.  The coordinate equations of \eqref{eq:system} read
$E^z_j=E^w_j=0$, where
\[
  E^z_j=z_j-\sum_{l\ne j}m_l\delta_{jl}^3(z_j-z_l),\qquad
  E^w_j=w_j-\sum_{l\ne j}m_l\delta_{jl}^3(w_j-w_l).
\]
Along a germ, $a\prec b$, $a\preccurlyeq b$ and $a\approx b$ in the sense of \cite[\S3]{AK} mean
$\ord a>\ord b$, $\ord a\ge\ord b$ and $\ord a=\ord b$ (Lemma~\ref{lem:germ}).

We prove that Diagram~1 germs do not exist (Theorem~\ref{thm:d1}).  The pair torque excludes a
triangle body at the $w$-centre of the triangle (Lemma~\ref{lem:d1-central}); then every germ
converges in a weighted chart (Proposition~\ref{prop:d1-cover}); its triangle tends to the
zero-multiplier triangle, which forces relation~(23) of \cite[\S8.1]{AK}
(Proposition~\ref{prop:d1-first}); and the potential of the triangle is too large for the
Morse--Bott estimate at that triangle.

\subsection{The chart}
Let $\gamma$ be a Diagram~1 germ at $m\in\Rp^5$; renumbering the bodies and exchanging $z$ and $w$
if necessary (symmetries of \eqref{eq:system}), its diagram is the one described above.  The
coordinate equations are homogeneous of
degree one in $z$ and in $w$ separately, and the ties involve only the products $z_{jl}w_{jl}$;
so $(z,w)\mapsto(cz,c^{-1}w)$ preserves them, and it does not change the diagram
\cite[p.~547]{AK}.
By Rule~2e \cite[p.~552]{AK} the three $z$-separations of the triangle have one order, and so do
its three $w$-separations.  Taking for $c$ a power of the parameter $t$ we assume that the
$w$-separations of the triangle have order $0$; the gauge $z_{12}=w_{12}$ is not used below.
Define $\lambda,\mu,\rho,\sigma$ by
\[
  \ord z_{12}=-\lambda,\qquad \ord z_{kl}=\mu,\qquad \ord Z_P=\mu+\rho,\qquad \ord W_P=\sigma,
\]
with $\rho=+\infty$ (resp.\ $\sigma=+\infty$) if $Z_P$ (resp.\ $W_P$) vanishes identically, in
which case the terms carrying $\hat Z$ (resp.\ $\hat W$) below are absent.  Here $\sigma$ is an
exponent, not the conjugation of Section~\ref{sec:ak}.
After a finite ramification of $t$ we write
\begin{equation}\label{eq:d1-chart}
\begin{gathered}
  z_i=t^{\mu+\rho}\hat Z+t^{-\lambda}A_i,\qquad w_i=t^{\sigma}\hat W+t^{\lambda}B_i
  \qquad(i\in P),\\
  z_k=-\tfrac{M_P}{M_T}\,t^{\mu+\rho}\hat Z+t^{\mu}\zeta_k,\qquad
  w_k=-\tfrac{M_P}{M_T}\,t^{\sigma}\hat W+\eta_k\qquad(k\in T),\\
  \delta_{12}=D_{12},\qquad\delta_{ik}=t^{\lambda/2}D_{ik},\qquad\delta_{kl}=t^{-\mu/2}D_{kl},
\end{gathered}
\end{equation}
where $(A_1,A_2)=(m_2,-m_1)\,\pi$, $(B_1,B_2)=(m_2,-m_1)\,\varpi$ and
$\sum_km_k\zeta_k=\sum_km_k\eta_k=0$.  The chart coordinates are $\hat Z,\hat W,\pi,\varpi$, the
mass-centred triangle offsets $\zeta=(\zeta_k)$, $\eta=(\eta_k)$, and the ten $D_{jl}$.  The
equations get the weights $t^{\lambda}$ for $E^z_i$, $t^{-\lambda}$ for $E^w_i$, $t^{\mu/2}$ for
$E^z_k$ and $t^{3\mu/2}$ for $E^w_k$; the ten ties keep weight $1$.
The exponents of interest form the open cone
\begin{equation}\label{eq:d1-cone}
  C=\{\lambda>0,\ \mu>0,\ \sigma-\lambda>0,\ \rho+\tfrac32\mu>0,\ \rho+\lambda+\mu>0\}.
\end{equation}
For the triangle block put, with sums over $T$,
\[
  \phi_k=\sum_{l\ne k}m_lD_{kl}^3(\zeta_k-\zeta_l),\qquad
  \psi_k=\sum_{l\ne k}m_lD_{kl}^3(\eta_k-\eta_l),\qquad
  \chi_{kl}=D_{kl}^2(\zeta_l-\zeta_k)(\eta_l-\eta_k)-1,
\]
$U_T=\sum_{k<l}m_km_lD_{kl}$ and $I_T=\sum_km_k\zeta_k\eta_k$; note
$\sum_km_k\phi_k=\sum_km_k\psi_k=0$ identically.  The crossed forces are
$X^z_k=\sum_{i\in P}m_i\delta_{ik}^3(z_k-z_i)$ and $X^w_k=\sum_{i\in P}m_i\delta_{ik}^3(w_k-w_i)$.

\begin{lemma}[computer-assisted]\label{lem:d1-chart}
Substitute \eqref{eq:d1-chart} into the ten coordinate equations and the ten ties and apply the
weights.  Each of the twenty results is a finite sum $\sum_ec_e\,t^e$, where the exponents $e$
are linear forms in $(\lambda,\mu,\rho,\sigma)$, $e=0$ occurs, and every $c_e$ is a polynomial in
the chart coordinates.  Every exponent $e\ne0$ that occurs is a combination of the five forms in
\eqref{eq:d1-cone} with nonnegative rational coefficients, at least one of them positive.  No
$c_0$ involves $\hat Z$ or $\hat W$.  For $k\in T$,
\[
  t^{\mu/2}E^z_k=t^{\mu/2}(z_k-X^z_k)-\phi_k,\qquad
  t^{3\mu/2}E^w_k=t^{3\mu/2}(w_k-X^w_k)-\psi_k,
\]
and the exponents occurring in $t^{\mu/2}(z_k-X^z_k)$, respectively $t^{3\mu/2}(w_k-X^w_k)$, are
\[
\begin{gathered}
  \tfrac32\mu,\ \tfrac12(\lambda+\mu),\ \tfrac32(\lambda+\mu),\ \tfrac32\mu+\rho,\
  \tfrac32(\lambda+\mu)+\rho,\\
  \text{respectively}\\
  \tfrac32\mu,\ \tfrac32(\lambda+\mu),\ \tfrac52\lambda+\tfrac32\mu,\ \sigma+\tfrac32\mu,\
  \tfrac32(\lambda+\mu)+\sigma .
\end{gathered}
\]
\end{lemma}
We carry out the substitution exactly, with the masses as indeterminates and the exponents as
linear forms in $(\lambda,\mu,\rho,\sigma)$.  The computation determines the exponents that
occur, writes each $e\ne0$ as a combination of the five forms with explicit nonnegative rational
coefficients, and checks the displayed identities for the triangle equations by expansion
(Appendix~\ref{app:verification}).

A direct computation gives the coefficients $c_0$.  They are $A_i(1-M_PD_{12}^3)$ and
$B_i(1-M_PD_{12}^3)$ for the equations of the pair, $-\phi_k$ and $-\psi_k$ for those of the
triangle, and $M_P^2D_{12}^2\pi\varpi-1$, $-D_{ik}^2A_i\eta_k-1$ and $\chi_{kl}$ for the ties of
the pair, of the six crossed pairs and of the triangle.
The \emph{boundary fibre} $\mathcal B_1$ is their common zero set in the coordinates other than
$\hat Z,\hat W$.  On $\mathcal B_1$ we have $D_{12}^3=M_P^{-1}$ and $D_{ik}^2A_i\eta_k=-1$, and the
triangle coordinates $(\zeta,\eta,D_{34},D_{35},D_{45})$ solve the \emph{zero-multiplier triangle
system} $\phi=\psi=0$, $\chi=0$.

\subsection{Coverage}
The coverage proof uses Rules 1e and 2e and Estimates 1 and 2 of \cite[\S3]{AK}
(Fact~\ref{fact:rules}) and \cite[Prop.~1]{AK} (Fact~\ref{fact:prop1}).

\begin{lemma}\label{lem:d1-syzygies}
The following are identities of polynomials:
\begin{gather*}
  m_1E^w_1+m_2E^w_2=m_1w_1+m_2w_2-\sum_{i\in P,\,k\in T}m_im_k\delta_{ik}^3(w_i-w_k),\\
  \sum_{k\in T}m_kE^z_k=\sum_{k\in T}m_kz_k-\sum_{i\in P,\,k\in T}m_im_k\delta_{ik}^3(z_k-z_i),\\
  T_P:=\sum_{i\in P,\,k\in T}m_im_k\delta_{ik}^3(z_iw_k-w_iz_k)
  =-\sum_{i\in P}m_i\big(w_iE^z_i-z_iE^w_i\big).
\end{gather*}
In particular the torque $T_P$ of the triangle on the pair vanishes on $A_m$.
\end{lemma}

\begin{proof}
In each sum the terms between two bodies of $P$, or between two bodies of $T$, cancel in pairs.
\end{proof}

\begin{lemma}\label{lem:d1-orders}
Let $\gamma$ be a Diagram~1 germ, normalized as above.
\begin{romanlist}
  \item $\ord\delta_{12}=0$ and $\ord w_{12}=\lambda>0$; the least order of the components of
  $\mathcal Z$ is $\ord z_1=\ord z_2=\ord Z_{12}=-\lambda$.
  \item $\mu>0$, $\ord\delta_{kl}=\ord Z_{kl}=-\mu/2$, and $\lambda>\mu/2$.
  \item $\ord z_{ik}=-\lambda$, $\ord w_{ik}<\lambda$ and $\ord w_i\ge\lambda$; hence
  $\ord W_T=\ord W_P\ge\lambda$.
  \item For $k\in T$ let $o_k$ be the coefficient of $t^0$ in $w_k-W_T$.  At most one $o_k$
  vanishes.  If $o_k\ne0$, then $\ord w_{ik}=0$ and $\ord\delta_{ik}=\lambda/2$.  If $o_k=0$, then
  $f=\ord(w_k-W_T)$ satisfies $0<f<\lambda$, and $\ord w_{ik}=f$, $\ord\delta_{ik}=(\lambda-f)/2$.
\end{romanlist}
\end{lemma}

\begin{proof}
(i) The isolated component $\{1,2\}$ of the $z$-diagram has circled vertices, so it has a
maximal $z$-stroke by Rule~1e \cite[p.~546]{AK}, which can only be $12$.  Estimate~2
\cite[p.~548]{AK} gives $r_{12}\approx1$, so $\ord\delta_{12}=0$ and $\ord w_{12}=-\ord
z_{12}=\lambda$.  Estimate~1 \cite[p.~548]{AK} with $z$ and $w$ exchanged gives
$w_{12}\approx\epsilon^2$, and Estimate~2 with $z$ and $w$ exchanged gives $w_{34}\succcurlyeq\epsilon$
for the $w$-stroke $34$.  Hence $w_{12}\prec w_{34}$; a comparison within one colour is not
affected by the normalizations, so $\lambda>\ord w_{34}=0$.  Finally $1,2$ are $z$-circled and
$12$ is a $z$-stroke, and $\ord Z_{12}=\ord z_{12}+3\ord\delta_{12}=-\lambda$.

(ii) Bodies $1,2$ carry no $w$-stroke, so $T$ is an isolated component of the $w$-diagram
without $w$-circles.  By Rule~1e it has no maximal $w$-stroke.  Estimate~2 with $z$ and $w$
exchanged gives $r_{kl}\preccurlyeq1$, with $r_{kl}\approx1$ only for a maximal $w$-edge, so
$r_{kl}\prec1$; thus $\mu=\ord(z_{kl}w_{kl})=2\ord r_{kl}>0$, $\ord\delta_{kl}=
-\mu/2$ and $\ord Z_{kl}=\mu-\frac32\mu=-\frac12\mu$.  Diagram~1 has no $z$-stroke inside $T$, so
$\ord Z_{kl}>-\lambda$, that is, $\lambda>\mu/2$.

(iii) The bodies of $T$ are not $z$-circled, so $\ord z_k>-\lambda=\ord z_i$ and $\ord
z_{ik}=-\lambda$.  Since $Z_{ik}^2=(z_{ik}w_{ik}^3)^{-1}$, $\ord Z_{ik}=\frac12(\lambda-3\ord
w_{ik})$, and the absence of a $z$-stroke $ik$ gives $\ord Z_{ik}>-\lambda$, i.e.\ $\ord
w_{ik}<\lambda$.  By Estimate~1 ($z$ and $w$ exchanged) $w_{2k}\succ\epsilon^2\approx w_{12}$,
because $12$ is the only maximal $z$-edge, and $z_{2k}\approx z_{12}$.  So
\cite[Prop.~1, p.~559]{AK} with $z$ and $w$ exchanged, applied to bodies $1,2$ as in the proof of
\cite[Prop.~4, p.~564]{AK}, gives $w_1,w_2\preccurlyeq w_{12}$, i.e.\ $\ord w_i\ge\lambda$.  This
also follows from Lemma~\ref{lem:d1-syzygies}: every term of $M_PW_P=\sum m_im_k\delta_{ik}^3
(w_i-w_k)$ has order $\frac12(3\lambda-\ord w_{ik})>\lambda$, and $w_i-W_P$ has order
$\ord w_{12}=\lambda$.

(iv) Since $\sum_km_ko_k=0$, and the $o_k$ do not all vanish because $\ord w_{kl}=0$, at most one
$o_k$ vanishes.  Now $w_k-w_i=(w_k-W_T)+(W_T-w_i)$ with $\ord(W_T-w_i)\ge\lambda>0$ by (iii).  If
$o_k\ne0$ this gives $\ord w_{ik}=0$.  If $o_k=0$ then $f>0$; $f\ge\lambda$ would give $\ord
w_{ik}\ge\lambda$, against (iii), so $f<\lambda$ and $\ord w_{ik}=f$.  In both cases
$\ord\delta_{ik}=-\frac12(\ord z_{ik}+\ord w_{ik})$.
\end{proof}

\begin{lemma}[computer-assisted]\label{lem:d1-torque}
Let $\lambda,\mu,f>0$ with $f<\lambda$ and let $\rho,\sigma$ be real.  Modify
\eqref{eq:d1-chart} by $(\eta_3,\eta_4,\eta_5)=(t^f\omega,\ m_5\vartheta,\
-m_3m_5^{-1}t^f\omega-m_4\vartheta)$ and $\delta_{i3}=t^{(\lambda-f)/2}d_i$ $(i=1,2)$, all other
entries unchanged.  Then
\[
  T_P=t^{(\lambda-f)/2}\Big(m_3\sum_{i=1,2}m_id_i\,\big(d_i^2A_i\omega\big)+\sum_{g\in G}c_g\,t^g\Big),
\]
where the $c_g$ are polynomials in the chart coordinates and $G$ consists of the thirteen
exponents
\[
\begin{gathered}
  \tfrac12f,\ \tfrac32f,\ \sigma-f,\ \sigma+\tfrac12f,\ 2\lambda+\mu-f,\
  2\lambda+\mu+\tfrac12f,\ \lambda+\mu+\rho,\ \lambda+\mu+\sigma-f,\ 2\lambda+\mu+\rho-f,\\
  \lambda+\mu+\rho+\tfrac12f,\ \lambda+\mu+\sigma+\tfrac12f,\ 2\lambda+\mu+\rho+\tfrac12f,\
  \lambda+\mu+\rho+\tfrac32f,
\end{gathered}
\]
each a combination of $\lambda,\ \mu,\ \sigma-\lambda,\ \rho+\lambda+\mu,\ f,\ \lambda-f$ with
nonnegative rational coefficients, at least one of them positive.  Moreover, if
$d_i^2A_i\omega=-1$ for $i=1,2$, then $(m_1d_1)^2-(m_2d_2)^2=-(m_1^3+m_2^3)/(m_1A_1\omega)$.
\end{lemma}
We expand $T_P$ exactly, with the masses as indeterminates; no relation among the masses is used,
in particular not (23).  Each exponent in $G$ is written as a combination of the six forms with
explicit nonnegative rational coefficients, and the last identity is checked by rational
simplification (Appendix~\ref{app:verification}).

\begin{lemma}\label{lem:d1-central}
Let $m\in\Rp^5$.  Along every Diagram~1 germ all three leading triangle $w$-offsets $o_k$ are
nonzero.
\end{lemma}

\begin{proof}
Suppose not.  By Lemma~\ref{lem:d1-orders}(iv), after renumbering $T$, $o_3=0\ne o_4,o_5$,
$0<f<\lambda$, $\ord\delta_{i3}=\frac12(\lambda-f)$ and $\ord\delta_{ik}=\frac12\lambda$ for $k=4,5$.
In the first two identities of Lemma~\ref{lem:d1-syzygies} the crossed terms have orders
$\frac12(3\lambda-\ord w_{ik})$ and $\frac12(\lambda-3\ord w_{ik})$, so
$\sigma\ge\frac12(3\lambda-f)$ and $\mu+\rho\ge\frac12(\lambda-3f)$, that is,
$\sigma-\lambda\ge\frac12(\lambda-f)>0$ and $\rho+\lambda+\mu\ge\frac32(\lambda-f)>0$.
Hence $\gamma$ is described by the chart of Lemma~\ref{lem:d1-torque}, with chart coordinates of
order $\ge0$, of order $0$ for $d_i,A_i,\omega$ (so these are nonzero at $t=0$), and every
exponent in $G$ is positive.
As $\ord(Z_T-Z_P)=\mu+\rho>-\lambda$ and $\sigma>\lambda>f$, we have $z_3-z_i=-t^{-\lambda}A_i(1+o(1))$
and $w_3-w_i=t^f\omega(1+o(1))$, so the ties $\delta_{i3}^2(z_3-z_i)(w_3-w_i)=1$ give
$d_i^2A_i\omega=-1$ at $t=0$.
By Lemma~\ref{lem:d1-torque}, $t^{-(\lambda-f)/2}T_P$ tends along $\gamma$ to
$-m_3(m_1d_1+m_2d_2)$, evaluated at $t=0$, and this is nonzero because
$(m_1d_1)^2-(m_2d_2)^2=-(m_1^3+m_2^3)/(m_1A_1\omega)\ne0$.  But $T_P=0$ along $\gamma$ by
Lemma~\ref{lem:d1-syzygies}.
\end{proof}

\begin{proposition}[Coverage]\label{prop:d1-cover}
Let $m\in\Rp^5$ and let $\gamma$ be a Diagram~1 germ, normalized as above.
\begin{romanlist}
  \item $\lambda>\mu/2>0$.
  \item Every leading triangle $w$-offset $o_k$ is nonzero; hence $\ord w_{ik}=0$ and
  $\ord\delta_{ik}=\lambda/2$ for $i\in P$, $k\in T$.
  \item $\sigma\ge\frac32\lambda$ and $\mu+\rho\ge\frac12\lambda$.  In particular
  $(\lambda,\mu,\rho,\sigma)\in C$, with $\sigma-\lambda\ge\frac12\lambda$,
  $\rho+\frac32\mu\ge\frac12(\lambda+\mu)$ and $\rho+\lambda+\mu\ge\frac32\lambda$.
  \item Along $\gamma$ every chart coordinate of \eqref{eq:d1-chart} converges as $t\to0$; the
  limits of $\pi$, $\varpi$, of the $D_{jl}$, of the $\eta_k$ and of the differences
  $\zeta_k-\zeta_l$, $\eta_k-\eta_l$ $(k\ne l)$ are nonzero; and the limit point lies in the
  boundary fibre $\mathcal B_1$.
\end{romanlist}
\end{proposition}

\begin{proof}
(i) is Lemma~\ref{lem:d1-orders}(i),(ii), and (ii) follows from Lemmas~\ref{lem:d1-central}
and~\ref{lem:d1-orders}(iv).  (iii) On $\gamma$, Lemma~\ref{lem:d1-syzygies} gives
$M_PW_P=\sum m_im_k\delta_{ik}^3(w_i-w_k)$ and $M_TZ_T=\sum m_im_k\delta_{ik}^3(z_k-z_i)$.  By (ii)
and Lemma~\ref{lem:d1-orders}(iii) the terms of the first sum have order $\frac32\lambda$ and
those of the second $\frac32\lambda-\lambda=\frac12\lambda$; as $\ord Z_T=\ord Z_P=\mu+\rho$ the
two bounds follow, and the three facet bounds are arithmetic.
(iv) By the definitions of $\lambda,\mu,\rho,\sigma$, by Lemma~\ref{lem:d1-orders} and by (ii),
every chart coordinate is a Puiseux series of order $\ge0$, of order exactly $0$ for the listed
ones; so all converge.  By Lemma~\ref{lem:d1-chart} and (iii), each weighted equation equals
$c_0+\sum_{e>0}c_et^e$ along $\gamma$ with bounded $c_e$; it vanishes identically, so $c_0$
vanishes at the limit point.
\end{proof}

\begin{remark}\label{rem:d1-ak}
(a) Lemma~\ref{lem:d1-orders} and Proposition~\ref{prop:d1-cover} give the clustering scheme
$z:1\ldots345\ldots2$, $w:3\ldots12\ldots5\ldots4$, which
\cite[\S8.1, p.~572]{AK} derives from Rule~2e, Estimate~2 and \cite[Prop.~4]{AK}.
Lemma~\ref{lem:d1-central} is the torque argument of the proof of \cite[Prop.~4]{AK}, carried out
in the chart with every competing term controlled.  The proof of Theorem~\ref{thm:d1} does not
use the clustering scheme of \cite[\S8.1]{AK}.
(b) If $o_3=0$, then $w_{34}:w_{45}:w_{53}=m_5:-(m_4+m_5):m_4$ at leading order.  Relation (21)
of \cite{AK}, which the triangle $w$-rows give at leading order, together with the ties and
$z_{34}+z_{45}+z_{53}=0$, then forces $m_3^2(m_4+m_5)^4=(m_4m_5)^3$.  On the sheet of (23) written, up to scale, as
$(m_3,m_4,m_5)=(b^2,b^2(b+1)^2,(b+1)^2)$, this relation is
$b^4(b-1)^2(b+1)^8(b^2+b+1)(b^4+b^3+4b^2+b+1)=0$, and with body $4$ or $5$ in place of body $3$ it
has no positive root.  So under (23) a vanishing $o_k$ can occur only on the central-body slice
$m_4=m_5=4m_3$.  We do not use this: Lemma~\ref{lem:d1-central} holds at every positive mass.
\end{remark}

\subsection{The mass condition}

\begin{lemma}[computer-assisted]\label{lem:d1-shape}
Let $m_3,m_4,m_5>0$ and let $(\zeta,\eta,D_{34},D_{35},D_{45})$, with
$\sum_km_k\zeta_k=\sum_km_k\eta_k=0$, solve the zero-multiplier triangle system.  Put
$\alpha_{kl}=\zeta_l-\zeta_k$ and $\beta_{kl}=\eta_l-\eta_k$ for $k<l$ in $T$.  Then
\begin{alphalist}
  \item $\beta_{kl}=\Lambda^{-1}\alpha_{kl}$ for one $\Lambda\in\C^*$, and $(D_{kl}\alpha_{kl})^2=
  \Lambda$ for all $k<l$;
  \item $(D_{34}\alpha_{34},D_{35}\alpha_{35},D_{45}\alpha_{45})=s\,(-1,1,-1)$ with $s^2=\Lambda$;
  the other six sign patterns admit no solution;
  \item $m_5\alpha_{34}^2=m_3\alpha_{45}^2=m_4\alpha_{35}^2$ and $Q(m_3,m_4,m_5)=0$, where
  \[
    Q=2m_3m_4m_5(m_3+m_4+m_5)-m_3^2m_4^2-m_4^2m_5^2-m_5^2m_3^2 .
  \]
\end{alphalist}
\end{lemma}
The proof is an exact computation with $m_3,m_4,m_5$ as indeterminates.  Explicit cofactor
identities give (a), and in each sign pattern the relation $\alpha_{35}=\alpha_{34}+\alpha_{45}$
eliminates the separations and leaves a polynomial relation among the masses: for the six
patterns excluded in (b) the polynomial is a square plus a polynomial with positive
coefficients, and for the pattern $s\,(-1,1,-1)$ it is $-Q$ (Appendix~\ref{app:verification}).
For positive masses $Q=0$ is equivalent to (23): with $x_k=m_k^{-1/2}$,
$Q=(m_3m_4m_5)^2(x_3+x_4+x_5)(-x_3+x_4+x_5)(x_3-x_4+x_5)(x_3+x_4-x_5)$, and the first factor is
positive.

The torus $\mathbb T=\{(\theta_1,\theta_2,\theta_3)\in(\C^*)^3:\theta_3^2=\theta_1\theta_2\}$ acts by
$(\zeta,\eta,D)\mapsto(\theta_1\zeta,\theta_2\eta,\theta_3^{-1}D)$.  It multiplies $\phi$ by
$\theta_1\theta_3^{-3}$, $\psi$ by $\theta_2\theta_3^{-3}$ and $I_T$ by $\theta_1\theta_2$, and it
fixes $\chi$.

\begin{lemma}[computer-assisted]\label{lem:d1-mb}
Let $b,c>0$, $(m_3,m_4,m_5)=(b^2c^2,\,b^2(b+c)^2,\,c^2(b+c)^2)$ and $\kappa>0$.  Let $y^*$ be the
point with $\eta=\zeta$, $\sum_km_k\zeta_k=0$, $\zeta_4-\zeta_3=\kappa m_5^{-1/2}$,
$\zeta_5-\zeta_3=-\kappa m_4^{-1/2}$, $D_{34}=m_5^{1/2}/\kappa$, $D_{35}=m_4^{1/2}/\kappa$ and
$D_{45}=-m_3^{1/2}/\kappa$.  In the coordinates $y=(\zeta_3,\zeta_4,\eta_3,\eta_4,D_{34},D_{35},
D_{45})$ let $F=(\phi_3,\phi_4,\psi_3,\psi_4,\chi_{34},\chi_{35},\chi_{45})$.  Then $F(y^*)=0$; the
Jacobian matrix $F'(y^*)$ has rank $5$, and its kernel is spanned by the generators
$(\zeta_3,\zeta_4,0,0,-\frac12D)$ and $(0,0,\eta_3,\eta_4,-\frac12D)$ of the $\mathbb T$-action; and
\[
  I_T(y^*)=\frac{\kappa^2}{m_3+m_4+m_5}\Big(\frac{m_3m_4}{m_5}+\frac{m_3m_5}{m_4}
  +\frac{m_4m_5}{m_3}\Big)>0 .
\]
\end{lemma}
The computation is exact in the field $\mathbb Q(b,c,\kappa)$: it gives $F(y^*)=0$, shows that
$F'(y^*)$ annihilates the two generators, and factors a $5\times5$ minor of $F'(y^*)$ and the
value $I_T(y^*)$ into factors with positive coefficients, so that both are positive for
$b,c,\kappa>0$ (Appendix~\ref{app:verification}).  One checks that every positive solution of
$m_3^{-1/2}=m_4^{-1/2}+m_5^{-1/2}$ is of the stated form.
By the equivariance above, the same rank and kernel statements hold at every point of
$O=\mathbb T\cdot y^*$, and $I_T\ne0$ on $O$.

\begin{proposition}[First condition]\label{prop:d1-first}
Let $m\in\Rp^5$ and let $\gamma$ be a Diagram~1 germ.  Then $Q(m_3,m_4,m_5)=0$, that is, the
triangle masses satisfy relation (23) of \cite[\S8.1]{AK}.  After renumbering $T$ so that
$m_3^{-1/2}=m_4^{-1/2}+m_5^{-1/2}$, the limit $y_0$ of
$y(t)=(\zeta_3,\zeta_4,\eta_3,\eta_4,D_{34},D_{35},D_{45})(t)$ lies on the orbit $O$ of
Lemma~\ref{lem:d1-mb}.
\end{proposition}

\begin{proof}
By Proposition~\ref{prop:d1-cover}(iv) the limit lies in $\mathcal B_1$, so its triangle part
solves the zero-multiplier triangle system.  In the notation of Lemma~\ref{lem:d1-shape},
$\phi_4=m_3\alpha_{34}D_{34}^3-m_5\alpha_{45}D_{45}^3$ and
$\phi_5=m_4\alpha_{45}D_{45}^3+m_3\alpha_{35}D_{35}^3$, and $m_3\phi_3=-m_4\phi_4-m_5\phi_5$;
likewise for $\psi$.  So $\phi=0$ and $\psi=0$ are relations (22) and (21) of \cite[\S8.1]{AK} in
cross-multiplied form, and Lemma~\ref{lem:d1-shape} gives $Q=0$.
After the renumbering, Lemma~\ref{lem:d1-shape}(c) and $\alpha_{35}=\alpha_{34}+\alpha_{45}$
leave only $(\alpha_{34},\alpha_{45},\alpha_{35})=\kappa(x_5,-x_3,-x_4)$ with $\kappa\in\C^*$, and
Lemma~\ref{lem:d1-shape}(a),(b) then determine $\eta$ and $D$ from $(\kappa,s)$.  The torus
$\mathbb T$ acts transitively on these pairs $(\kappa,s)$, and $y^*$ is the solution with
$\kappa>0$, $s=-1$.  Hence the solutions form the single orbit $O$, and $y_0\in O$.
\end{proof}

\subsection{Exclusion}

\begin{lemma}[computer-assisted]\label{lem:d1-virial}
For arbitrary $\zeta,\eta\in\C^3$ and $D_{34},D_{35},D_{45}\in\C$,
\begin{gather*}
  \sum_{k\in T}m_k\zeta_k\psi_k=U_T+\sum_{k<l}m_km_lD_{kl}\chi_{kl},\\
  2\,dU_T+\sum_{k\in T}m_k\big(\psi_k\,d\zeta_k+\phi_k\,d\eta_k\big)
  =\sum_{k<l}m_km_l\big(D_{kl}\,d\chi_{kl}-2\chi_{kl}\,dD_{kl}\big).
\end{gather*}
\end{lemma}
Both identities are checked by exact expansion, with the masses as indeterminates
(Appendix~\ref{app:verification}).
On the tie manifold $\{\chi=0\}$ they reduce to $U_T=\sum_km_k\zeta_k\psi_k$ and
$dU_T=-\frac12\sum_km_k(\psi_k\,d\zeta_k+\phi_k\,d\eta_k)$; both vanish on $O$.

\begin{lemma}\label{lem:d1-estimate}
Every $y_0\in O$ has a neighbourhood $V\subset\C^7$ with a constant $c_V$ such that
$|U_T(y)|\le c_V\,|(\phi_3,\phi_4,\psi_3,\psi_4)(y)|^2$ for every $y\in V$ with $\chi(y)=0$.
\end{lemma}

\begin{proof}
By Lemma~\ref{lem:d1-mb} and equivariance, five components $F_5$ of $F$ have independent
differentials at $y_0$.  The orbit map $\theta\mapsto\theta y_0$ is an immersion of a
neighbourhood of $1\in\mathbb T$ onto a $2$-dimensional piece of $O$ inside the $2$-dimensional
manifold $F_5^{-1}(0)$ near $y_0$, so $F_5^{-1}(0)\cap V\subset O$ for small $V$, and
$\operatorname{dist}(y,O)\le c\,|F_5(y)|\le c\,|(\phi_3,\phi_4,\psi_3,\psi_4)(y)|$ for $y\in V$
with $\chi(y)=0$ (implicit function theorem).  Near $y_0$ the manifold $\{\chi=0\}$ is the graph
of $D$ over $(\zeta_3,\zeta_4,\eta_3,\eta_4)$, since $\partial\chi_{kl}/\partial D_{kl}=
2D_{kl}^{-1}(\chi_{kl}+1)\ne0$.  In these coordinates $U_T$ and its differential vanish on $O$
(Lemma~\ref{lem:d1-virial}), so Taylor's formula at a nearest point of $O$ gives $|U_T(y)|\le
c'\operatorname{dist}(y,O)^2$.
\end{proof}

\begin{theorem}\label{thm:d1}
Let $m\in\Rp^5$.  No curve in $A_m$ has a Diagram~1 germ (Definition~\ref{def:germ}).  In
particular there is no Diagram~1 germ with an unbounded four-product, which is the case $r=1$ of
Theorem~\ref{thm:exclusion}\,(a).
\end{theorem}

\begin{proof}
Let $\gamma$ be a Diagram~1 germ, normalized as above and with $T$ renumbered as in
Proposition~\ref{prop:d1-first}.  Along $\gamma$ the ties of the triangle give $\chi(y(t))=0$
exactly, and $y(t)\to y_0\in O$.

\emph{Forcing.}  Along $\gamma$ we have $E^z_k=E^w_k=0$, so Lemma~\ref{lem:d1-chart} gives
$\phi_k=t^{\mu/2}(z_k-X^z_k)$ and $\psi_k=t^{3\mu/2}(w_k-X^w_k)$, with the exponents listed
there and bounded coefficients.  By Proposition~\ref{prop:d1-cover}(iii),
$\frac32\mu+\rho\ge\frac12(\lambda+\mu)$.  Hence $|\phi(y(t))|=O(|t|^r)$ with
$r=\min(\frac32\mu,\frac12(\lambda+\mu))>0$, and $|\psi(y(t))|=O(|t|^{3\mu/2})$.

\emph{Morse--Bott.}  By Lemma~\ref{lem:d1-estimate}, $|U_T(y(t))|=O(|t|^{2r})$.

\emph{Virial.}  By Lemma~\ref{lem:d1-virial} and $\chi=0$,
$U_T(y(t))=\sum_km_k\zeta_k\psi_k=t^{3\mu/2}\sum_km_k\zeta_k(w_k-X^w_k)$.  Since
$w_k=W_T+\eta_k$ and $\sum_km_k\zeta_k=0$, we get $\sum_km_k\zeta_kw_k=I_T(y(t))$, while
$\ord X^w_k\ge\frac32\lambda$ by Proposition~\ref{prop:d1-cover}(ii).  Hence
$U_T(y(t))=t^{3\mu/2}\big(I_T(y_0)+o(1)\big)$ with $I_T(y_0)\ne0$.

\emph{Contradiction.}  The last two steps give $\frac32\mu\ge2r=\min(3\mu,\lambda+\mu)$.  Since
$\mu>0$, this means $\lambda+\mu\le\frac32\mu$, i.e.\ $\lambda\le\frac12\mu$, which contradicts
Proposition~\ref{prop:d1-cover}(i).
\end{proof}

\begin{remark}\label{rem:d1-wall}
The window $\lambda>\mu/2$ is used only in the last step.  In the chart family
\eqref{eq:d1-chart} the outer bodies are $z$-circled exactly when $\lambda\ge\mu/2$.  The wall
$\lambda=\mu/2$ is the Diagram~4 chart after a rotation; there the orders $\frac32\mu$ and
$\lambda+\mu$ coincide and the argument gives nothing, which is why Diagram~4 needs the analysis
of Section~\ref{sec:d45}.  For $\lambda<\mu/2$ the diagram is Diagram~5.  No bound on
four-products is used, so Theorem~\ref{thm:d1} excludes Diagram~1 germs as first and as second
limits.
\end{remark}
